La propagación electromagnética en un medio con pérdidas se estudió mediante la evolución de los campos y el transporte de energía. Para una antena inmersa, la disminución de la señal dependió de la expansión geométrica, la absorción del material y la distribución espacial del campo. El vector de Poynting proporcionó una relación entre estas magnitudes y la conservación de energía \citep{Griffiths2017,Poynting1884}.

En este trabajo se implementó una simulación en Python con el motor FDTD Meep a 3 GHz. Se representó un monopolo en coordenadas cilíndricas axisimétricas y se consideraron cuatro medios homogéneos: vacío, concreto, agua salada y tejido muscular. La frecuencia situó el estudio en el intervalo sub-6 GHz; no se simuló una red 5G, su modulación ni un enlace completo entre transmisor y receptor.

El análisis reunió dos conjuntos de evidencia. Diez configuraciones permitieron comparar perfiles radiales de potencia y estimar una atenuación aparente. Cinco corridas adicionales evaluaron el balance entre flujo entrante, flujo saliente y disipación de Joule en un volumen anular sin fuentes. La comparación principal se realizó a 5 px/mm. Las resoluciones adicionales permitieron examinar la aproximación a la teoría y la sensibilidad a la malla.

Se distinguieron la incertidumbre condicional del ajuste, la discrepancia respecto de la teoría, la variación entre discretizaciones y el residuo energético. Esta separación evitó interpretar un intervalo estrecho o un buen cierre del balance como demostración suficiente de exactitud. Los códigos, los registros y las figuras se organizaron para conservar la trazabilidad de cada resultado.

El alcance se limitó a la propagación dentro de cada medio, con geometrías escaladas mediante su longitud de onda. No se midió transmisión a través de paredes o interfaces, ni se estimó la eficiencia global de una antena calibrada. La discusión evaluó el cumplimiento de los objetivos dentro de estas condiciones y delimitó las verificaciones necesarias para extender el modelo.
